\subsection{Generalized influence for all exchangeable laws}
\label{ian:exchangeable-generalized-influence}

Status: proved. The conclusion covers all exchangeable laws, including
endpoint slices, with constants uniform under arbitrary
positive-probability pinnings. The product influence input is
Corollary~\ref{ian:product-influence-aa-proved}, proved in
Subsection~\ref{ian:influence-form-uniform-aa} from (BC) and the OSSS
covariance inequality; the slice step cites Girish
Proposition~4.12(1) and Lemma~4.9. An earlier version assumed the imported uniform
theorem (UAA) of
Section~\ref{ht:conventions-and-imported-uniform-theorem}; that
hypothesis is no longer used. The recomputed constants are recorded in
Proposition~\ref{ian:influence-consequences-exchangeable}.

Notation \(D_i\), \(I_i^\mu\), \(G_\mu\) is that of
Subsection~\ref{ian:generalized-influence-bridge}.

Here are explicit, deliberately coarse constants:
\[
 \kappa=c_\star=\frac{9}{2^{61}\pi^{12}},\qquad
 \lambda=10^{-9}2^{-146}\kappa^4,\qquad
 \sigma=2^{-112}\lambda,
 \tag{5}\label{ian:exchangeable-constants}
\]
where \(c_\star\) is the constant of
Theorem~\ref{ian:uniform-influence-aa-proved}.

Let \(\rho_{n,p}\) be the product law with
\(\Pr(X_i=1)=p\), and let \(\omega_{n,k}\) be uniform on the slice with
exactly \(k\) coordinates equal to \(1\).

\begin{lemma}[Product and slice bounds uniform in density]
\label{ian:uniform-product-slice-aa}
The following statements hold for every \(n,d\ge1\),
every \(t\in(0,1]\), and every globally bounded multilinear
\(f:\{-1,1\}^n\to[-1,1]\) of degree at most \(d\):
\[
 \Var_{\rho_{n,p}}f\ge t
 \quad\Longrightarrow\quad
 \max_i I_i^{\rho_{n,p}}(f)\ge\kappa(t/d)^{12}
 \quad\text{for every }p\in[0,1],
 \tag{6}\label{ian:uniform-product-bound}
\]
and
\[
 \Var_{\omega_{n,k}}f\ge t
 \quad\Longrightarrow\quad
 \max_i I_i^{\omega_{n,k}}(f)\ge\lambda(t/d)^{56}
 \quad\text{for every }k\in\{0,\ldots,n\}.
 \tag{7}\label{ian:uniform-slice-bound}
\]
\end{lemma}

\begin{proof}
The bound \eqref{ian:uniform-product-bound} is (P-AA) of
Corollary~\ref{ian:product-influence-aa-proved}, proved in
Subsection~\ref{ian:influence-form-uniform-aa}, applied to the product
law \(\rho_{n,p}\); that corollary holds for every bias, including
\(p\in\{0,1\}\), with no bias-dependent constant. If
\(\Var_{\rho_{n,p}}f\ge t\), it gives
\[
 \max_i I_i^{\rho_{n,p}}(f)
 \ge c_\star\frac{(\Var_{\rho_{n,p}}f)^{12}}{d^{12}}
 \ge\kappa(t/d)^{12}.
\]
See also
Section~\ref{influence-recovery:a-product-aa-bound-uniform-over-all-biases},
which derives the better rate \(s^{7/2}/d^7\), with inexplicit
constants, from its own (UAA) hypothesis by dyadic encoding; that
route is not used here.

For the slice bound we use Girish Proposition~4.12(1), pp.~20--21,
whose conventions are \(\{0,1\}\), range \([0,1]\) and unnormalised
flip influence. Under the translation of
Subsection~\ref{ian:program}, \(g(y)=(1+f(2y-1))/2\) has
\(\Var(g)=\Var(f)/4\) and
\(\operatorname{Inf}^{\mathrm{Girish}}_i(g)=I_i(f)\). Thus
\eqref{ian:uniform-product-bound} supplies the Girish product bound
\(S(u,d)=\kappa(u/d)^{12}\), uniformly in \(p\). For \(0<k<n\), apply
Proposition~4.12 at \(p=k/n\), where \(\nu_{n,k/n}=\omega_{n,k}\) since \(\lfloor(k/n)n\rfloor=k\), with
slice variance parameter \(t/4\). Its bound contains no \(p\): the proof
uses only Lemma~2.12 (constant \(1/(2n^2)\)) and Corollary~4.11
(constant \(48d^2/\sqrt n\)). It gives
\[
 \max_i I_i^{\omega_{n,k}}(f)
 \ge10^{-9}d^{-8}
 \min\left\{(t/4)^5,\
       (t/4)\kappa^4\left(\frac{t}{8d}\right)^{48}\right\}
 =10^{-9}2^{-146}\kappa^4\,t^{49}d^{-56}
 \ge\lambda(t/d)^{56},
\]
the minimum being its second term since \(\kappa,t\le1\le d\). At
\(k=0,n\) the variance vanishes, so \eqref{ian:uniform-slice-bound} is
vacuous.
\end{proof}

\begin{lemma}[Slice continuity]
\label{ian:slice-continuity}
For \(f\) as above, define
\[
 F(k)=\E_{\omega_{n,k}}f,\quad
 H(k)=\E_{\omega_{n,k}}f^2,\quad
 V(k)=\Var_{\omega_{n,k}}f,\quad
 J_i(k)=I_i^{\omega_{n,k}}(f).
\]
For every \(k,u\in\{0,\ldots,n\}\),
\[
 \begin{split}
 |F(k)-F(u)|&\le8d^2|k-u|/n,\\
 |H(k)-H(u)|&\le16d^2|k-u|/n,\\
 |V(k)-V(u)|&\le32d^2|k-u|/n,\\
 |J_i(k)-J_i(u)|&\le16d^2|k-u|/n.
 \end{split}
 \tag{10}\label{ian:slice-lipschitz-estimates}
\]
In particular,
\[
 V(k)\le32d^2\min\{k/n,1-k/n\}.
 \tag{11}\label{ian:endpoint-variance}
\]
\end{lemma}

\begin{proof}
Apply Girish Lemma~4.9, p.~19, respectively to
\((1+f)/2\), \(f^2\), and \((D_i f)^2\).
The latter two have range \([0,1]\) and multilinear degree at most
\(2d\). These applications give the bounds for \(F,H,J_i\).
Since \(|F|\le1\) and \(V=H-F^2\), the bound for \(V\) follows.
Finally \(V(0)=V(n)=0\), which proves
\eqref{ian:endpoint-variance}.
\end{proof}

\begin{theorem}[Theorem B]
\label{ian:exchangeable-gamma}
The class \(\Lambda_{\mathrm{exch}}\) of all
exchangeable laws on all finite cubes is pinning-closed and satisfies
\(\Gamma^{\mathrm{rand}}(\Lambda_{\mathrm{exch}})\). One choice of
constants is
\[
 c_L=\frac{\sigma^2}{264},\qquad C_L=114,\qquad
 c_G=\frac1{18},\qquad C_G=1,\qquad
 B_G=128,\qquad K_G=4,
 \tag{12}\label{ian:exchangeable-gamma-constants}
\]
where \(\sigma\) is defined in
\eqref{ian:exchangeable-constants}. Consequently, by the observation after
Definition~\ref{ian:gamma-definition},
\(\Gamma^{\mathrm{set}}(\Lambda_{\mathrm{exch}})\) also holds.
\end{theorem}

\begin{proof}
Pinning closure: fix a positive-probability assignment on some
coordinates. Every permutation of the remaining coordinates extends to
one fixing the pinned coordinates; it preserves the exchangeable law and
the conditioning event, hence the residual law. Nothing below depends on
the law beyond exchangeability, so the constants are uniform over all
pinnings.

Fix \(n,d\ge1\), an exchangeable \(\mu\), \(v\in(0,1]\), and a
globally bounded degree-at-most-\(d\) polynomial \(f\) with
\(\Var_\mu f\ge v\). Let
\[
 W=\#\{i:X_i=1\},\qquad
 \pi_k=\Pr_\mu(W=k).
\]
Exchangeability gives
\(\mu=\sum_{k=0}^n\pi_k\omega_{n,k}\). With the notation of
Lemma~\ref{ian:slice-continuity}, total variance gives
\[
 \Var_\mu f=\Var(F(W))+\E V(W).
 \tag{13}\label{ian:exchangeable-variance-split}
\]
Set \(\ell=\min\{n,\lceil128(d/v)^4\rceil\}\).
If \(\ell=n\), the Global alternative holds with gain
\(\Var_\mu f\ge v\). Hence assume \(\ell<n\).

\emph{Between slices.} Suppose first that \(s:=\Var(F(W))\ge v/2\).
Let \(S\) be an independent uniformly random \(\ell\)-subset, and let
\(\widetilde F\) be the piecewise linear interpolation on \([0,1]\) of
the values \(F(k)\) at \(k/n\); by
\eqref{ian:slice-lipschitz-estimates} it is \(8d^2\)-Lipschitz. Define
\[
 \widehat p=\frac1\ell\sum_{i\in S}\frac{1+X_i}{2},
 \qquad Z=\widetilde F(\widehat p).
\]
All expectations in this part are over \((S,X)\). Hypergeometric
sampling gives
\(\E[\widehat p\mid X]=W/n\) and
\(\E[(\widehat p-W/n)^2\mid X]\le1/(4\ell)\), so
\[
 \delta:=\E(Z-F(W))^2
 \le\frac{16d^4}{\ell}
 \le\frac{v^4}{8}\le\frac v8\le\frac s4.
 \tag{14}\label{ian:density-estimation-error}
\]
For fixed \(X=x\), \(\E_S Z\) depends only on the weight of \(x\);
this is where the random \(S\) is used. Since
\(\E[f-F(W)\mid W]=0\), with \(m=\E_\mu f\),
\[
 \E[(f-m)(Z-m)]
 =\E[(F(W)-m)(Z-m)]
 \ge s-\sqrt{s\delta}\ge s/2,
 \tag{15}\label{ian:between-slice-covariance}
\]
while \(\E(Z-m)^2\le(\sqrt s+\sqrt\delta)^2\le9s/4\). As \(Z\) is
measurable with respect to \(\mathcal H=\sigma(S,X_S)\),
Cauchy--Schwarz gives
\[
 \E_S G_\mu(f,S)
 =\E\bigl(\E[f\mid\mathcal H]-m\bigr)^2
 \ge
 \frac{\bigl(\E[(f-m)(Z-m)]\bigr)^2}{\E(Z-m)^2}
 \ge \frac s9\ge\frac v{18}\ge\frac1{18}(v/d).
\]
This is the Global alternative.

\emph{Within slices.} Otherwise \(\E V(W)\ge v/2\) by
\eqref{ian:exchangeable-variance-split}. Let
\(\mathcal A=\{k:V(k)\ge v/4\}\). Because \(0\le V(k)\le1\),
\[
 v/2\le\E V(W)\le v/4+\pi(\mathcal A),
 \qquad\text{so}\qquad \pi(\mathcal A)\ge v/4.
 \tag{16}\label{ian:good-slice-mass}
\]
For every \(k\in\mathcal A\),
Lemma~\ref{ian:uniform-product-slice-aa} with \(t=v/4\) gives a
coordinate with influence at least
\[
 s_*:=\lambda\left(\frac{v}{4d}\right)^{56}
      =\sigma(v/d)^{56}.
\]
The influential coordinate may vary with \(k\), so we localise in
density. Partition \([0,1]\) into
\(q=\lceil32d^2/s_*\rceil\le33d^2/s_*\) intervals of length \(1/q\).
By \eqref{ian:good-slice-mass}, one interval contains points \(k/n\)
with \(k\in\mathcal A\) of total \(\pi\)-mass at least \(v/(4q)\).
Choose such a \(k_0\) with \(\pi_{k_0}>0\) and a coordinate \(i\) with
\(J_i(k_0)\ge s_*\). By \eqref{ian:slice-lipschitz-estimates}, every
\(k\) in that interval has
\(J_i(k)\ge s_*-16d^2/q\ge s_*/2\). Hence, using \(v\le1\),
\[
 I_i^\mu(f)
 =\sum_k\pi_kJ_i(k)
 \ge\frac{s_*}{2}\frac{v}{4q}
 \ge\frac{vs_*^2}{264d^2}
 \ge\frac{\sigma^2}{264}(v/d)^{114}.
\]
This is the Local alternative.
\end{proof}

\begin{remark}[Endpoints, margins, and the simulation theorem]
\label{ian:exchangeable-endpoints}
There is no endpoint exclusion in Theorem~\ref{ian:exchangeable-gamma}.
By \eqref{ian:endpoint-variance}, a slice with \(V(k)\ge v/4\) has
\(k/n\in[v/(128d^2),\,1-v/(128d^2)]\), so the within-slice argument
has a polynomial density margin automatically. Since
Lemma~\ref{ian:uniform-product-slice-aa} is uniform in density, a
pinning that moves the residual density towards \(0\) or \(1\) needs no
separate treatment; constants for a fixed margin \([a,1-a]\) would not
suffice.

Combined with Theorem~\ref{ian:bridge}, this gives hereditary full AA
for all exchangeable laws, with the influence input proved in
Subsection~\ref{ian:influence-form-uniform-aa}. The acceptance-simulation
conclusion is already established in
\path{simulation/exchangeable_inputs/01_theorem.tex},
Section~\ref{exch:heading-2}, from a simulation-form uniform
hypothesis, which is supplied by
Corollary~\ref{cor:resolution-product-alphabet-aa} applied to the
uniform law; an adaptive-transcript influence form with better
exponents, conditional on (UAA), is
Section~\ref{influence-recovery:the-theorem}. The
additional conclusion here is the generalized-influence dichotomy, with
a random fixed-size set, for every globally bounded low-degree
polynomial.
\end{remark}
